\documentclass{article}
\usepackage[T1]{fontenc}
\usepackage{lmodern}
\usepackage{graphicx}
\usepackage{amsmath,amssymb}
\usepackage[a4paper,margin=2cm]{geometry}
\usepackage[absolute,overlay]{textpos}
\setlength{\TPHorizModule}{1mm}
\setlength{\TPVertModule}{1mm}
\usepackage[indonesian]{babel}
\usepackage{hyperref}
\setlength{\emergencystretch}{2em}
% unit: Halaman 31

\begin{document}

% === Halaman 31 (llm) ===

\begin{itemize}
  \item Karakteristik \textit{cipher}alir adalah melokalisir kesalahan enkripsi/dekripsi hanya pada bit yang eror saja.
  \item Maksudnya, kesalahan 1-bit pada cipherteks hanya menghasilkan kesalahan 1-bit pada plainteks hasil dekripsi.
  \item Misalkan cipherteks $1010\, \mathbf{0}010\, 1000\, 1000\, 1000\, 1111$ (Hex: A2888F) mengalami eror pada bit ke-5 dari semula 0 menjadi 1: $1010\, \mathbf{1}010\, 1000\, 1000\, 1000\, 1111$
  \item Ketika didekripsi:
\end{itemize}

\vspace{5mm}

\begin{align*}
\text{C: } & 1010\mathbf{1}0101000100010001111... \\
\text{K: } & 111101011001000111101011... \oplus \\
\hline
& 0101\mathbf{1}1110001100101100100... \quad (5\mathbf{F}1964\text{ dalam Hexa}) \\
& \quad \quad \quad \quad \quad \quad \quad \quad \text{(Seharusnya: 571964 dalam Hexa)}
\end{align*}

\end{document}
